% =====================================================================
% BARIS "Model prediksi" TERTAHAN.
%
% Dua alasan lama - keterjelasan kontribusi fitur dan kebutuhan
% komputasi rendah - SUDAH DICABUT oleh KEPUTUSAN #5:
%   - keterjelasan  : digugurkan J7. Yang dimiliki RF secara bawaan
%                     adalah MDI, sedangkan MDI tidak boleh dilaporkan karena
%                     bias kardinalitas. Angka yang dilaporkan adalah
%                     permutation importance, yang MODEL-AGNOSTIK.
%   - komputasi     : terbalik secara faktual (RF paling mahal).
%
% Isi baris tersebut setelah keputusan Random Forest atau Gradient
% Boosting ditetapkan. Rujukan: docs/KEPUTUSAN_DIAMBIL.md,
% docs/ALUR_SISTEM_FINAL.md, docs/REVISI_BAB_II_4_1.md.
% =====================================================================
\begin{longtable}{@{}p{2.5cm}p{2.7cm}p{3.8cm}p{3.6cm}@{}}
\caption{Alasan pemilihan dan kelemahan yang diakui.}
\label{tab:penentuan-solusi} \\
\toprule
\textbf{Keputusan} & \textbf{Alternatif terpilih} & \textbf{Alasan pemilihan} & \textbf{Kelemahan yang diakui} \\
\midrule
\endfirsthead

\multicolumn{4}{@{}l}{\small\emph{Tabel \thetable{} (lanjutan)}} \\
\toprule
\textbf{Keputusan} & \textbf{Alternatif terpilih} & \textbf{Alasan pemilihan} & \textbf{Kelemahan yang diakui} \\
\midrule
\endhead

\bottomrule
\endlastfoot

Model prediksi & \emph{Gradient Boosting} & Kriteria pemilihannya akurasi kelas kondisi, bukan galat regresi, sebab kelas kondisi itulah yang menyusun kueri penelusuran; pada kriteria tersebut ia unggul, sementara galat regresinya setara & Kepekaan terhadap hipoglikemia pada horizon 60~menit lebih rendah daripada \emph{Random Forest}, dan pengklasifikasinya menukar ketepatan demi kepekaan \\[2pt]

Model bahasa & Layanan awan & Tidak membebani perangkat lokal dan tersedia berjenjang gratis & Memerlukan konektivitas jaringan sebagaimana Batasan~5 \\[2pt]

Model \emph{embedding} & \emph{sentence-transformers} lokal & Berjalan tanpa akselerator grafis dan tanpa biaya per permintaan & Keterbatasan pada penelusuran lintas bahasa sebagaimana Subbab~\ref{subsec:arsitektur-rag} \\[2pt]

Kerangka antarmuka & Streamlit & Pengembangan cepat untuk purwarupa dan integrasi langsung dengan ekosistem Python & Keterbatasan penyesuaian tata letak untuk kebutuhan lanjutan \\[2pt]

Skema pengambilan pengetahuan & RAG di atas korpus tak terstruktur & Tidak menuntut pembangunan ontologi maupun integrasi rekam medis & Mutu bergantung pada cakupan korpus \\[2pt]

Pengklasifikasi kondisi & Pengklasifikasi terpisah & Terbukti menangkap kejadian hipoglikemia jauh lebih baik daripada penetapan ambang atas nilai regresi & Menambah satu model yang perlu dilatih dan dipelihara, serta menurunkan nilai duga positif \\[2pt]

Orkestrasi alur & Penulisan sebagai kode program & Memenuhi pemisahan logika domain (KNF-07) dan berjalan tanpa server khusus (KNF-09) & Menuntut penulisan kode yang lebih banyak dibandingkan perkakas visual \\

\end{longtable}
